\begin{abstract}
Non-convulsive seizures affect approximately 30\% of critically ill patients who undergo EEG
monitoring, yet optimal treatment strategies remain uncertain because of the trade-off between
seizure suppression and drug toxicity. Randomized trials in this setting are expensive, ethically
fraught, and have so far been inconclusive, while the wide practice variation in routine ICU care
creates observational data in which treatment decisions and disease progression influence each other
over time---a confounding structure that defeats standard survival analysis.

We present CICADAS (\textbf{C}ausal \textbf{I}nference for \textbf{C}ritical-Care
\textbf{A}nti-Seizure Treatment with \textbf{D}isease Dynamics, \textbf{A}utomated PKPD/Trial
Simulation, and \textbf{S}urvival Optimization), a computational framework that couples the
parametric g-formula with mechanistic pharmacokinetic--pharmacodynamic modeling, disease dynamics,
and dynamic treatment policies to emulate randomized trials from longitudinal observational data. We
specify the notation and identification assumptions required for time-varying treatment and
confounding, describe a sequential estimation procedure for each model component, and provide open
implementations in MATLAB and Python.

The framework is evaluated entirely by \emph{internal, simulation-based validation} on synthetic ICU
cohorts with known ground truth. Across 400 simulated replications, confounding by indication
attenuated a true 14.6 percentage-point survival benefit to 1.0 points under naive Kaplan--Meier
analysis---eliminating 93\% of the effect---and reversed its sign in 39\% of replications. Under
these conditions the g-formula recovers the ground-truth treatment effect closely, whereas naive
Kaplan--Meier, stabilized inverse probability of treatment weighting, and a baseline-weighted
marginal structural model do not. We report the estimator's
behavior under injected unmeasured confounding, EEG measurement error, and alternative censoring
mechanisms.

These results establish methodological feasibility, not clinical utility. Because the simulated data
and the estimation models share structural assumptions, the recovery results characterize estimator
consistency under correct specification rather than robustness to misspecification. Whether CICADAS
yields valid causal estimates in real ICU data is an open question, to be addressed by external
validation in multicenter ICU EEG cohorts.
\end{abstract}
